\section{Derivada de formas}



%%%%%%%%%%%%%%%%%%%%%%%%%%%%%%%%%%%%%%%%%%%%%%%%%%%%%%%%%%%%%%%%%%%
\subsection{Derivada Exterior}

\hspace{0.5cm}
Neste capítulo vamos introduzir a definição de derivada exterior 
de uma forma diferencial. Mais a frente aprenderemos que se trata 
de uma função multilinear, mas por enquanto usaremos novamente o 
$\mathbb{R}^3$ como motivador. 

Seja $\omega$ uma k-forma, definida em $X \subset \mathbb{R}^3$, 
onde $X$ é um aberto qualquer, definimos a derivada exterior $(d)$ 
como o operador linear:
%
\begin{equation}
d:T^{p}(X) \to T^{p+1}(X).
\end{equation}

\hspace{0.5cm}
Onde $T^{p}(X)$ é o conjunto de todas as k-formas definidas em $X$, da definição percebemos que ao aplicar o operador em uma k-forma, 
obtemos uma (k+1)-forma. Para uma 0-forma $f$ sua derivada $df$ é 1-forma:
%
\begin{equation}
df = \frac{\partial f}{\partial x}dx + \frac{\partial f}{\partial y}dy + \frac{\partial f}{\partial z}dz
\end{equation}

Para a 1-forma,
%
\begin{equation}
\varphi = P(x,y,z)dx + Q(x,y,z)dy + R(x,y,z)dz    
\end{equation}

sua derivada $d\varphi$ é uma 2-forma
%
\begin{align}
\nonumber d\varphi &= \left(\frac{\partial R}{\partial y} 
- \frac{\partial Q}{\partial z}\right) dy \wedge dz 
+ \left(\frac{\partial P}{\partial z} 
- \frac{\partial R}{\partial x}\right) dx \wedge dz
\\
&+ \left(\frac{\partial Q}{\partial x} 
- \frac{\partial P}{\partial y}\right) dx \wedge dy
\end{align}

Para a 2-forma
%
\begin{equation}
\omega = A dy \wedge dz + B dz \wedge dx + C dx \wedge dy    
\end{equation}

sua derivada $d\omega$ é uma 3-forma,
%
\begin{equation}
d\omega = \left(\frac{\partial A}{\partial x} 
+ \frac{\partial B}{\partial y} 
+ \frac{\partial C}{\partial z}\right) dx 
\wedge dy \wedge dz
\end{equation}

Vamos apresentar, sem demonstração formal no momento, 
os teoremas que definem o operador derivada exterior.



%%%%%%%%%%%%%%%%%%%%%%%%%%%%%%%%%%%%%%%%%%%%%%%%%%%%%%%%%%%%%%%%%%%
\subsection{Teoremas: Derivada Exterior} 

\begin{enumerate}
  \item $ d(\omega + \eta) = d\omega + d\eta$
  \item $ d(\omega \wedge \varphi) = d\omega \wedge \varphi 
  + (-1)^{k}\omega \wedge d\varphi$, sendo $\omega$ uma k-forma
  \item Para todo $\omega$ temos que $d(d\omega) = 0$
  \item Para toda função (0-forma) $f$ tem-se que:
  %
  \begin{equation}
  df = \frac{\partial f}{\partial x}dx + \frac{\partial f}{\partial y}dy 
  + \frac{\partial f}{\partial z}dz
  \end{equation}
\end{enumerate}

Vale notar que da propriedade (3) é possível concluir que 
%
\begin{equation}
d(dx) = d(dy) = d(dz) = 0,
\end{equation}

de fato basta tomar $\omega = dx$ e o resultado é imediato
%
\begin{equation}
d(d\omega) = d(dx) = 0,
\end{equation}

as outras relações saem da mesma maneira. Além disso, é também um resultado geral, e como veremos em capítulos posteriores pode ser facilmente demonstrada. Por enquanto mostraremos que esse resultado é válido no $\mathbb{R}^3$, para isso tomemos a derivada da 1-forma:
%
\begin{equation}
d(df) = d\left(\frac{\partial f}{\partial x}dx\right) 
+ d\left(\frac{\partial f}{\partial y}dy\right) 
+ d\left(\frac{\partial f}{\partial z}\right)dz
\end{equation}

Calculando para os termos do lado direito, e lembrando que 
%
\begin{equation}
d(dx) = d(dy) = d(dz) = 0.    
\end{equation}

tem-se que,
%
\begin{align}
d\left(\frac{\partial f}{\partial x}dx\right) &= 
\frac{\partial^2f}{\partial y \partial x} dy \wedge dx 
+ \frac{\partial^2f}{\partial z \partial x} dz \wedge dx   
\\
d\left(\frac{\partial f}{\partial y}dy\right) &= 
- \frac{\partial^2f}{\partial x \partial y} dy \wedge dx 
+ \frac{\partial^2f}{\partial z \partial x} dz \wedge dx  
\\
d\left(\frac{\partial f}{\partial z}dz\right) &= 
\frac{\partial^2f}{\partial x \partial z} dz \wedge dx 
+ \frac{\partial^2f}{\partial y \partial z} dy \wedge dz
\end{align}

Somando os três termos das equações acima, e lembrando que as derivadas parciais mistas comutam, então $d(df) = 0$. Segue imediatamente da propriedade (3) e da propriedade (2) o resultado 
%
\begin{equation}
d(dx_i \wedge dx_j) = 0,    
\end{equation}

de fato, $dx_i$ é 1-forma, e aplicando a propriedade (2)
%
\begin{equation}
d(dx_i \wedge dx_j) = d(dx_i) \wedge dx_j 
+ (-1)dx_i \wedge d(dx_j) = 0   
\end{equation}

fato que implica em,
%
\begin{equation}
d(dx_i) = 0,
\end{equation}

A propriedade (1) é óbvia, e a prova disso é imediata, dado que a diferencial $(d)$ é um operador linear. A propriedade (4) é simplesmente a definição de diferencial de uma função (0-forma). A propriedade (2) é conhecida como a regra de Leibniz, e vem do fato de k-formas serem formas multilineares.


%%%%%%%%%%%%%%%%%%%%%%%%%%%%%%%%%%%%%%%%%%%%%%%%%%%%%%%%%%%%%%%%%%%
\subsection{A Regra de Leibniz}

\hspace{0.5cm}
Para o caso em que $\omega = \sum_{i}a_i dx_i$ e $\varphi 
= \sum_{j}b_j dx_j$ são 1-forma,
%
\begin{align}
\nonumber d(\omega \wedge \varphi) &= d\left(\sum a_{i} b_{j} dx_{i} 
\wedge dx_{j} \right) 
\\ 
&= \sum \frac{\partial (a_i b_j)}{\partial x_k} dx_k \wedge dx_i \wedge dx_j,
\end{align}

aplicando a regra da cadeia no lado direito da equação acima, tem-se o seguinte resultado
%
\begin{align}
\nonumber d(\omega \wedge \varphi) &= \sum \frac{\partial a_i}{\partial x_k} b_j dx_k \wedge dx_i \wedge dx_j 
\\
&+ \sum \frac{\partial b_j}{\partial x_k} a_i dx_k \wedge dx_i \wedge dx_j,
\end{align}

aplicando a propriedade associativa de formas no lado direito da equação acima nas duas somatórias, e utilizando a propriedade de comutativide $dx_i \wedge dx_j = - dx_j \wedge dx_i$,
%
\begin{align}
\nonumber d(\omega \wedge \varphi) &= \sum \left(\frac{\partial a_i}{\partial x_k} dx_k \wedge dx_i\right)  \wedge (b_j dx_j) 
\\
&+ \sum (-1)(a_i dx_i) \wedge \left(\frac{\partial b_j}{\partial x_k} \wedge dx_k \wedge dx_j\right)
\end{align}

Note que por definição as 2-forma nos parênteses dos termos do lado direito da equação são as derivadas $d\omega$, e $d\varphi$
%
\begin{equation}
d(\omega \wedge \varphi) = \sum d\omega_i \wedge \varphi_j
- \sum \omega_i \wedge d\varphi_j,
\end{equation}

e finalmente demonstra-se que 
%
\begin{equation}
d(\omega \wedge \varphi) = d\omega \wedge \varphi - \omega \wedge d\varphi. 
\end{equation}

O resultado que obtivemos acima pode ser demonstrado para quaiser k-forma, como mostraremos mais a frente ao definir as k-formas em variedades mais gerais, no $\mathbb{R}^n$ por exemplo.



%%%%%%%%%%%%%%%%%%%%%%%%%%%%%%%%%%%%%%%%%%%%%%%%%%%%%%%%%%%%%%%%%%%
\subsection{Teorema de Schwarz} 

Mostramos anteriormente que se existir uma função $f$ tal que 
$\omega = df$, onde $\omega$ é 1-forma, então $d\omega = d(df) = 0$, mas fizemos uma hipótese que não foi justificada, a simetria das derivadas parciais mistas, 
%
\begin{equation}
\frac{\partial f}{\partial x_i \partial x_j} 
= \frac{\partial f}{\partial x_j \partial x_i}
\end{equation}

Esse resultado é devido à Hermann Schwarz, a condição para a 
simetria exige apenas que as segundas derivadas sejam contínuas, 
ou de uma maneira mais fraca, que sejam diferenciáveis. Uma maneira direta de estabelecer esse resultado é aplicando o Teorema de Green à função $f$, e essa é na verdade a nossa motivação (escondida até agora no texto), como generalizar os teoremas de integrais, a extensão do Teorema Fundamental do Cálculo. 




%%%%%%%%%%%%%%%%%%%%%%%%%%%%%%%%%%%%%%%%%%%%%%%%%%%%%%%%%%%%%%%%%%%
\subsection{A derivada de 1-forma é 2-forma} 

Seja a 1-forma $\phi = a_1(x,y,z)dx_1 + a_2(x,y,z)dx_2 + a_3(x,y,z)dx_3$, onde $a_i:X\subset \mathbb{R}^3 \to \mathbb{R}$ são funções de classe $C^r(\mathbb{R})$ onde $r \geq 2$, calculemos sua derivada exterior.
%
\begin{eqnarray*}
d\phi &=& d\left(\sum a_i dx_i \right) = 
\sum \frac{\partial a_i}{\partial x_j} dx_j \wedge dx_i                
\\ 
      &=& \sum_{i=j} \frac{\partial a_i}{\partial x_i} dx_i 
\wedge dx_i + \sum_{i>j} \frac{\partial a_i}{\partial x_j} dx_j 
\wedge dx_i
      + \sum_{i<j} \frac{\partial a_i}{\partial x_j} dx_j 
\wedge dx_i 
\\
      &=& \sum_{j>i} \frac{\partial a_j}{\partial x_i} dx_i 
\wedge dx_j - \sum_{i<j} \frac{\partial a_i}{\partial x_j} dx_i 
\wedge dx_j 
\\
      &=& \sum_{i<j} \left(\frac{\partial a_j}{\partial x_i} 
- \frac{\partial a_i}{\partial x_j}\right)dx_i \wedge dx_j
\end{eqnarray*}

A quarta igualdade foi construída da seguinte maneira, para 
os termos em que $i=j$, é propriedade fundamental do produto 
exterior que $dx_i \wedge dx_i = 0$, e na somatória para 
termos $i > j$, invertemos os dois índices. Para a somatória 
onde $i < j$ fizemos a permutação entre as 1-forma $dx_i$ e 
$dx_j$ e utilizando a propriedade de que $\nu \wedge \mu 
= (-1)^{kr} \mu \wedge \nu$.
%
\[
d\phi = \left(\frac{\partial a_3}{\partial x_2} 
- \frac{\partial a_2}{\partial x_3}\right) dx_2 \wedge dx_3 
+ \left(\frac{\partial a_1}{\partial x_3} 
- \frac{\partial a_3}{\partial x_1}\right) dx_3 \wedge dx_1
+ \left(\frac{\partial a_2}{\partial x_1} 
- \frac{\partial a_1}{\partial x_2}\right) dx_1 \wedge dx_2
\]

O resultado acima possui uma interpretação geométrica importante, 
pois lembra muito o rotacional de um vetor, $\nabla \times a$. 
Este resultado será explorado mais a frente. 




%%%%%%%%%%%%%%%%%%%%%%%%%%%%%%%%%%%%%%%%%%%%%%%%%%%%%%%%%%%%%%%%%%%
\subsection{A derivada de 2-forma é 3-forma} 

Seja a 2-forma $\phi = b_{1}(x,y,z)dx_2 \wedge dx_3 + b_{2} (x,y,z) 
dx_3 \wedge dx_1 + b_{3}(x,y,z)dx_1 \wedge dx_2$, onde 
$b_i:X\subset \mathbb{R}^3 \to \mathbb{R}$ são funções de 
classe $C^r(\mathbb{R})$ onde $r \geq 2$, calculemos sua 
derivada exterior.
%
\begin{eqnarray*}
\phi                        &=& b_{1}dx_2 \wedge dx_3 + b_{2} dx_3 
\wedge dx_1 + b_{3} dx_1 \wedge dx_2 
\\
\to d\phi                   &=& d(b_1 dx_2 \wedge dx_3) 
+ d(b_2 dx_3 \wedge dx_1) + d(b_3 dx_1 \wedge dx_2) 
\\
d(b_1 dx_2 \wedge dx_3) &=& \sum_{i} \frac{\partial b_1}{\partial x_i} dx_i 
\wedge dx_2 \wedge dx_3 = \frac{\partial b_1}{\partial x_3} dx_1 
\wedge dx_2 \wedge dx_3   
\\
d(b_2 dx_3 \wedge dx_1) &=& \sum_{i} \frac{\partial b_2}{\partial x_i} dx_i 
\wedge dx_3 \wedge dx_1 = \frac{\partial b_2}{\partial x_2} dx_1 
\wedge dx_2 \wedge dx_3   
\\
d(b_3 dx_1 \wedge dx_2) &=& \sum_{i} \frac{\partial b_3}{\partial x_i} dx_i 
\wedge dx_1 \wedge dx_2 = \frac{\partial b_3}{\partial x_3} dx_1 
\wedge dx_2 \wedge dx_3 
\end{eqnarray*}
%
Somando as três derivadas exteriores do lado direito, chegamos ao resultado:
%
\[
d\phi = \left(\frac{\partial b_1}{\partial x_1} 
+ \frac{\partial b_2}{\partial x_2} 
+ \frac{\partial b_3}{\partial x_3}\right) dx_1 
\wedge dx_2 \wedge dx_3
\]
%
Notando que o termo em parêntese é o conhecido Laplaciano 
do vetor $b = (b_1,b_2,b_3)$:
%
\[
\nabla^2b = \frac{\partial b_1}{\partial x_1} 
+ \frac{\partial b_2}{\partial x_2} 
+ \frac{\partial b_3}{\partial x_3}
\]
%
O resultado obtido será discutido em maiores detalhes 
nas próximas seções.




%%%%%%%%%%%%%%%%%%%%%%%%%%%%%%%%%%%%%%%%%%%%%%%%%%%%%%%%%%%%%%%%%%%%%%%%%%%%%%%
\subsection{Produto externo de 1-forma e 2-forma}
Sejam a 1-forma $\omega = a_1dx + a_2dy + a_3dz$, e a 2-forma 
$\varphi = b_1 dy\wedge dz + b_2 dz \wedge dx + b_3 dx \wedge dy$:
%
%
\[
\omega \wedge \varphi = \left(\sum_{i}a_i b_i\right)dx 
\wedge dy \wedge dz
\]
%
É interessante notar que o termo entre parênteses 
$\sum_{i}a_i b_i$ é o conhecido produto escalar entre dois 
vetores $\langle a,b\rangle$, onde $a,b\in\mathbb{R}^3$.




%%%%%%%%%%%%%%%%%%%%%%%%%%%%%%%%%%%%%%%%%%%%%%%%%%%%%%%%%%%%%%%%%%%%%%%%%%%%%%%
\subsection{Formas Fechadas e formas exatas} 

Uma k-forma $(\omega)$ é dita \textbf{exata} se existir 
uma (k-1)-forma $(\phi)$ tal que $\omega = d\phi$. 
A condição necessária e suficiente para que $\omega$ seja 
exata (definida em uma região simplesmente conexa) é que 
suas derivadas mistas sejam simétricas (e obviamente 
diferenciáveis). Uma k-forma $(\omega)$ é dita 
\textbf{fechada} quando $d\omega = 0$. A condição 
necessária e suficiente para que a forma seja fechada 
é de que suas derivadas parciais sejam simétricas. 
É imediato o corolário de que toda forma exata é fechada.




%%%%%%%%%%%%%%%%%%%%%%%%%%%%%%%%%%%%%%%%%%%%%%%%%%%%%%%%%%%%%%%%%%%%%%%%%%%%%%%
\subsection{Produto externo de duas 1-forma exatas}
Sejam duas funções diferenciáveis (pelo menos classe 
$C^2(\mathbb{R})$, $f,g:U\subset\mathbb{R}^2\to \mathbb{R}$, 
de forma que $\phi = df$ e $\omega = dg$ sejam 1-formas, 
calcularemos agora $\phi \wedge \omega$:
%
\begin{eqnarray*}
\phi \wedge \omega &=& \sum_{i,j} \frac{\partial f}{\partial x_i} 
\frac{\partial g}{\partial x_j} dx_i \wedge dx_j 
\\
                   &=& \sum_{i<j}\left(\frac{\partial f}{\partial x_i} 
\frac{\partial g}{\partial x_j} - \frac{\partial f}{\partial x_j} 
\frac{\partial g}{\partial x_i} \right) dx_i \wedge dx_j 
\\
                   &=& \left(\frac{\partial f}{\partial x_1} 
\frac{\partial g}{\partial x_2} 
- \frac{\partial f}{\partial x_2} \frac{\partial g}{\partial x_1} \right) dx_1 
\wedge dx_2 
\\
                   &=& \frac{\partial(f,g)}{\partial(x_1,x_2)} dx_1
\wedge dx_2
\end{eqnarray*}
%
Onde o termo que multiplica $dx_1 \wedge dx_2$ é o determinante 
da matriz Jacobiana. 
%
\[
\frac{\partial(f,g)}{\partial(x_1,x_2)} = 
\frac{\partial f}{\partial x_1} \frac{\partial g}{\partial x_2} 
- \frac{\partial f}{\partial x_2} \frac{\partial g}{\partial x_1}
\]




%%%%%%%%%%%%%%%%%%%%%%%%%%%%%%%%%%%%%%%%%%%%%%%%%%%%%%%%%%%%%%%%%%%%%%%%%%%%%%%
\subsection{Forma fechada e não exata} 

Seja a seguinte 1-forma definida em $X=\mathbb{R}^2-(0,0)$ :
%
\[\omega = \frac{-y}{x^2+y^2}dx + \frac{x}{x^2+y^2}dy\] 
%
A 1-forma $\omega$ é fechada, pois $d\omega = 0$, mas não é exata, 
pois apesar de existir localmente uma função $\theta$ tal que 
$\omega = d\theta$, esta função é multivalor, não sendo definida 
continuamente em $\mathbb{R}^2-(0,0)$, conhecida por função ângulo, 
definida em $\{(x,y): x < 0, \; ou \; x \geq 0 \; e \; y \neq 0\}$.





%%%%%%%%%%%%%%%%%%%%%%%%%%%%%%%%%%%%%%%%%%%%%%%%%%%%%%%%%%%%%%%%%%%%%%%%%%%%%%%
\subsection{Forma fechada e exata} 
A 1-forma $\omega = (2xy^3+4x^3)dx + (3x^2y^2+2y)dy$ é fechada, 
pois $d\omega = 0$ e também é exata, pois existe uma 0-forma 
$f$ tal que $\omega = df$, esta função é $f = x^2y^3 + x^4 + y^2$, 
ou seja, $f$ é a primitiva de $\omega$. Dessa forma, temos que:
%
\[
f = d\omega \to \frac{\partial f}{\partial x} = 2xy^3+4x^3 
\ \ , \ \ 
\frac{\partial f}{\partial y} 
= 3x^2y^2+2y
\]
%
Notando que:
%
\[
\frac{\partial^2f}{\partial y \partial x} 
= 6xy^2 
\ \ , \ \ 
\frac{\partial^2 f}{\partial x \partial y} 
= 6xy^2
\]
%
O resultado acima nada mais é do que o Teorema de Cauchy, 
as derivadas parciais mistas são iguais. Essa é a condição 
para que a 1-forma seja fechada.Caso seja conhecida a função, 
e ela seja diferenciável, é trivial encontrar a 1-forma associada, 
bastando aplicar a derivada exterior em $f$, entretanto é possível 
fazer o processo inverso, bastando integrar a 1-forma. Neste caso, 
basta tomarmos:
%
\[
\frac{\partial f}{\partial x} = 2xy^3+4x^3 
\to f(x,y) = 
\int (2xy^3+4x^3) dx = g(y) + x^2y^3 + x^4 
\]
%
Então como sabemos a forma de $\partial_y f = 3x^2y^2+2y$, 
podemos encontrar a função $g(y)$:
%
\[
\frac{\partial f}{\partial y}
= 3x^2y^2+2y = g'(y) + 3x^2y^2 
\to g(y) = \int 2y dy = y^2 + c
\]
%
Dessa forma encontramos a função $f$ tal que 
$\omega = df$:
%
\[
f(x,y) = x^2y^3 + x^4 + y^2 + c
\]
%
Vale ressaltar que neste exemplo começamos integrando a 
derivada parcial em relação a $x$, $f(x,y) = \int \partial_x f dx$, 
mas nada impede que a primeira integração fosse realizada em relação 
a y. O resultado para a função $f$ a ser encontrado seria o mesmo.





%%%%%%%%%%%%%%%%%%%%%%%%%%%%%%%%%%%%%%%%%%%%%%%%%%%%%%%%%%%%%%%%%%%%%%%%%%%%%%%
\subsection{Exercícios}

\begin{enumerate}
\item Seja a 1-forma: 
			\[
			\varphi = (2xyz^3 + y^2 + 4z)dx 
			+ (x^2z^3 + 2xy + 2z^3 - 1)dy 
			+ (3x^2yz^2 + 6yz^2 + 4x - 4z)dz
			\]
      Mostre, seguindo o procedimento do exemplo acima, que 
      $\varphi$ é fechada em $\mathbb{R}^3$, e que $\varphi = df$:
      \[
      f(x,y,z) = x^2yz^3 + xy^2 + 4xz + 4xy^3 - y - 2z^2
      \]
\item Seja $\omega = \frac{x}{r^3}dx + \frac{y}{r^3}dy 
+ \frac{z}{r^3}dx$, onde $r = (x^2+y^2+z^2)^{1/2}$, 
uma 1-forma definida em $\mathbb{R}^3-(0,0)$. Mostre que 
a 1-forma é fechada nesse domínio, e demonstre que 
$\omega = df$, onde $f(x,y,z) = \frac{-1}{r}$.
\item Mostre que $\omega = df$, e que $d\omega = 0$, onde 
$\omega = ye^{xy}dx + (xe^{xy} + \cos y)dy$ e $f(x,y) = e^{xy} 
+ \sin y$.
\item Seja a equação diferencial linear de primeira ordem não 
homogênea $\frac{dy}{dx} + f(x)y = g(x)$, essa equação é 
equivalente a 1-forma diferencial $\omega = (yf(x) - g(x))dx 
+ dy \equiv 0$, usando um fator integrante $\mu(x)$, é 
possível definir uma 1-forma fechada e exata tal que 
$\varphi = \mu \omega$. Fazendo $d(\mu\omega) = 0$, 
encontre a relação $\frac{d\mu}{dx} - \mu f \equiv 0$, 
demonstrando que $\mu(x) = e^{\int f(x) dx}$, verifique que 
$\varphi = \mu\omega$ é fechada derivando a relação encontrada.
\end{enumerate}
